% !TEX program = xelatex
% ============================================================
% 编译方法：
%   cd /workspace/files && xelatex -interaction=nonstopmode tech_report.tex && xelatex -interaction=nonstopmode tech_report.tex && cd /workspace
% 输出: tech_report.pdf
% ============================================================
\documentclass[11pt,a4paper]{article}

% -- 中文支持 --
\usepackage{fontspec}
\usepackage{xeCJK}
\setCJKmainfont{AR PL SungtiL GB}[
  BoldFont=AR PL KaitiM GB,
]
\setCJKmonofont{AR PL SungtiL GB}

% -- 页面设置 --
\usepackage[top=2.5cm, bottom=2.0cm, left=2.5cm, right=2.5cm, headheight=14pt]{geometry}
\usepackage{fancyhdr}
\pagestyle{fancy}
\fancyhf{}
\fancyhead[C]{\small 海信 × 阿里云 AI 智能体大赛 --- 技术方案报告}
\fancyfoot[C]{\thepage}
\renewcommand{\headrulewidth}{0.4pt}

% -- 常用包 --
\usepackage{graphicx}
\usepackage{hyperref}
\usepackage{enumitem}
\usepackage{booktabs}
\usepackage{array}
\usepackage{xcolor}
\usepackage{listings}
\usepackage{titlesec}
\usepackage{caption}
\usepackage{float}

% -- 代码样式 --
\lstset{
  basicstyle=\small\ttfamily,
  backgroundcolor=\color{gray!5},
  frame=single,
  rulecolor=\color{gray!30},
  breaklines=true,
  columns=flexible,
  numbers=none,
  aboveskip=6pt,
  belowskip=6pt,
}

% -- 章节格式 --
\titleformat{\section}{\Large\bfseries}{\thesection}{1em}{}
\titleformat{\subsection}{\large\bfseries}{\thesubsection}{1em}{}
\titleformat{\subsubsection}{\normalsize\bfseries}{\thesubsubsection}{1em}{}
\titlespacing*{\section}{0pt}{12pt}{6pt}
\titlespacing*{\subsection}{0pt}{8pt}{4pt}
\titlespacing*{\subsubsection}{0pt}{6pt}{3pt}

% -- 正文 --
\begin{document}

% ==================== 封面 ====================
\begin{titlepage}
  \centering
  \vspace*{3cm}

  {\Huge\bfseries AI 智能体大赛 \\[6pt]
   技术方案报告}

  \vspace{1.5cm}

  {\LARGE AstraDraft \\[8pt]
   CAD 智能问答 Agent 系统}

  \vspace{3cm}

  {\large
   \begin{tabular}{rl}
     参赛队伍： & 中午吃什么 \\
     赛题名称： & 海信 × 阿里云 AI 智能体大赛 \\
     提交日期： & 2026 年 6 月 \\
   \end{tabular}
  }

  \vfill
\end{titlepage}

% ==================== 目录 ====================
\newpage
\tableofcontents
\newpage

% ==================== 1. 项目背景与目标 ====================
\section{项目背景与目标}

\subsection{赛题简介}

本赛题基于海信企业真实生产流程，以工业 CAD 图纸为核心数据载体，要求参赛者构建一个具备自然语言理解能力的智能 Agent 系统，能够深度解析 CAD 图纸，并通过自然语言问答的方式，自动检索并准确反馈图纸中的关键技术参数。

\subsection{核心任务}

\begin{enumerate}[nosep]
  \item \textbf{搭建智能 Agent 系统}：构建能够深度解析 CAD 图纸（DWG/DXF 格式）的系统，提取几何实体、尺寸标注、文字注释等信息。
  \item \textbf{精准参数提取}：支持用户通过自然语言提问（如“过滤网的高度是多少？”），系统自动检索并反馈图纸中的关键技术参数。
  \item \textbf{多图纸支持}：适应训练图纸（过滤网）以及考核图纸（侧板、印刷版）的解析和问答。
\end{enumerate}

\subsection{技术指标}

\begin{itemize}[nosep]
  \item 解析 CAD 图纸，构建包含 102+ 个结构化参数的语义索引。
  \item 支持 12 类以上查询意图（尺寸查询、计数、极值、比较、材质、性能等）。
  \item 训练集准确率 \textbf{93.1\%}（58 题通过评估）。
  \item 混合策略：规则管道优先，LLM 回退兜底（DeepSeek v4 Pro）。
\end{itemize}

% ==================== 2. 系统总体架构 ====================
\section{系统总体架构}

\subsection{架构总览}

系统采用分层架构设计，自底向上分为数据层、CAD 理解引擎、Agent 核心调度层和用户界面层，各层职责明确、接口清晰。

\begin{lstlisting}[language={},frame=single,title=系统架构图]
+-----------------------------------------+
|         用户界面层 (CLI)                 |
+----------------+------------------------+
                 | 自然语言问题
                 v
+-----------------------------------------+
|     Agent 核心调度层                      |
|  +----------+  +----------------------+  |
|  | 意图识别  |  | 规则管道 + LLM 回退  |  |
|  +----------+  +----------------------+  |
+----------------+------------------------+
                 |
                 v
+-----------------------------------------+
|       CAD 理解引擎                       |
|  +------------+  +------------------+   |
|  | 参数索引    |  | 语义映射(jieba)  |   |
|  +------------+  +------------------+   |
+----------------+------------------------+
                 |
+-----------------------------------------+
|       CAD 解析层                         |
|  DWG→DXF → ezdxf → CADDocument          |
|  (实体 / 尺寸 / 文本 / 图层)             |
+-----------------------------------------+
\end{lstlisting}

\subsection{分层职责}

\begin{table}[H]
\centering
\small
\begin{tabular}{llp{7cm}}
\toprule
层次 & 技术选型 & 职责 \\
\midrule
用户界面层 & Python CLI & 接收用户自然语言输入，展示回答 \\
Agent 核心调度层 & Python + LLM API & 意图识别、上下文管理、规则管道优先匹配、LLM 回退次选 \\
CAD 理解引擎 & ezdxf + jieba + 规则引擎 & 解析图纸、构建参数索引、语义别名映射 \\
数据层 & JSON / 文件系统 & CAD 文件存储、参数索引缓存 \\
\bottomrule
\end{tabular}
\end{table}

\subsection{数据流}

\begin{enumerate}[nosep]
  \item 输入 DWG → \texttt{dwg2dxf} 转换为 DXF 格式。
  \item ezdxf 读取 DXF，构建 \texttt{CADDocument}（含实体、标注、文本、图层）。
  \item 参数索引构建器从维度和文本中提取参数，建立结构化索引。
  \item 用户提问 → 意图分类 → 规则管道逐级尝试匹配 → 若均未命中则触发 LLM 回退。
  \item 格式化答案并输出。
\end{enumerate}

% ==================== 3. CAD 解析模块 ====================
\section{CAD 解析模块}

\subsection{模块功能}

CAD 解析模块（\texttt{cad\_parser/}）负责将 DWG/DXF 格式的 CAD 图纸文件转换为程序可以理解的结构化数据。它是整个系统的基础层，为上层语义理解提供原始数据。

\subsection{技术选型}

\begin{itemize}[nosep]
  \item \textbf{DWG → DXF 转换}：使用开源库 \textbf{LibreDWG} 的 \texttt{dwg2dxf} 工具，编译部署于 Docker 环境，支持 AutoCAD 2018（AC1032）格式。
  \item \textbf{DXF 解析}：使用 Python 库 \textbf{ezdxf}（$\ge$1.4.0）读取 DXF 文件，提取实体、尺寸标注和文本。
  \item \textbf{数据建模}：基于 \textbf{Pydantic v2} 定义严格的数据结构，确保类型安全。
\end{itemize}

\subsection{核心数据结构}

系统定义的核心数据模型（\texttt{cad\_parser/models.py}）：

\begin{itemize}[nosep]
  \item \textbf{CADDocument}：图纸文档的顶层表示，包含实体列表、尺寸列表、文本注释列表、图层信息、图块定义和元数据。
  \item \textbf{Entity}：几何实体（LINE、ARC、CIRCLE、LWPOLYLINE 等 14 种），包含图层、类型和几何参数。
  \item \textbf{Dimension}：尺寸标注（线性、对齐、直径、半径、角度共 7 种），包含测量值、标注文本、定义点坐标。
  \item \textbf{TextAnnotation}：文本注释，包含文字内容、位置、字高、旋转角度。
  \item \textbf{LayerInfo}：图层定义，包含颜色、线型、实体数量。
\end{itemize}

\subsection{文件结构}

\begin{lstlisting}[language={},frame=single]
cad_parser/
+-- __init__.py          # 模块入口
+-- reader.py            # DXF 文件读取器
+-- entities.py          # 几何实体提取
+-- dimensions.py        # 尺寸标注提取与解析
+-- text_extractor.py    # TEXT/MTEXT 文本提取
+-- layer_manager.py     # 图层管理
+-- models.py            # Pydantic 数据模型
\end{lstlisting}

\subsection{关键技术细节}

\begin{itemize}[nosep]
  \item \textbf{编码处理}：中国工业 CAD 图纸内部使用 GBK 编码，解析时正确解码中文标注文本。
  \item \textbf{单位处理}：当 \texttt{insunits=0}（未指定单位）时，按行业惯例默认使用毫米（mm）。
  \item \textbf{DimType 位标志}：ezdxf 返回的 \texttt{dimtype} 使用复合位标志，需通过位运算（\texttt{\& 0x0F}）提取真实标注类型。
  \item 训练图纸解析结果：3206 个实体、299 个文本注释、全部尺寸标注完整解析。
\end{itemize}

% ==================== 4. 语义理解与索引构建 ====================
\section{语义理解与索引构建}

\subsection{模块功能}

CAD 语义理解模块（\texttt{cad\_understanding/}）位于解析层之上，负责将原始的 CAD 实体和标注转化为结构化的参数索引，并建立自然语言术语与 CAD 参数之间的语义映射关系。

\subsection{参数索引构建流程}

参数索引的构建分为四个阶段：

\begin{enumerate}[nosep]
  \item \textbf{维度提取}：遍历所有尺寸标注，提取测量值（\texttt{measurement}），结合附近文字推断参数名称和语义别名。
  \item \textbf{键值对提取}：正则匹配标注文本中的 \texttt{键=值} 模式（如 \texttt{高度=245mm}），直接解析为参数。
  \item \textbf{领域启发式}：针对中国工业 CAD 图纸的特点，提取标题栏信息（图纸名称、图号、比例、图幅、分类、版本号）、签字栏信息（设计、审核、工艺、标准化、批准）、技术要求、BOM 表等结构化信息。
  \item \textbf{语义别名分配}：基于预定义的同义词映射表（\texttt{SYNONYM\_MAP}），为每个参数分配自然语言别名（如“高度$\rightarrow$[高, H, height]”）。
\end{enumerate}

\subsection{语义映射}

\texttt{semantic\_mapper.py} 使用 \textbf{jieba 分词} 结合多层级评分匹配：

\begin{itemize}[nosep]
  \item 精确匹配：参数名或别名直接匹配用户查询（得分 10）。
  \item 分词匹配：对用户查询分词后逐词比对（得分 5）。
  \item 主名称偏序：查询词出现在参数主名称中额外加分（得分 5）。
  \item 子串匹配：参数名或别名包含在词中（得分 2）。
  \item 同义词加权：预定义跨域同义词组命中（如"圆角$\rightarrow$半径"、"每英寸$\rightarrow$目数"、"树脂$\rightarrow$材质"，得分 12）。
\end{itemize}

\subsection{文件结构}

\begin{lstlisting}[language={},frame=single]
cad_understanding/
+-- __init__.py               # 模块入口
+-- parameter_index.py        # 参数索引构建（核心）
+-- semantic_mapper.py        # 语义映射（jieba + 同义词）
+-- component_recognizer.py   # 部件/特征识别
+-- scale_inferencer.py       # 单位推断
+-- models.py                 # 语义模型定义
\end{lstlisting}

\subsection{已构建的参数类别}

训练图纸共构建 108+ 个结构化参数，覆盖以下类别：

\begin{itemize}[nosep]
  \item 线性/对齐尺寸参数（长度、宽度、高度、间距等）
  \item 孔径/半径参数
  \item 标题栏信息（图纸名称、图号、比例、图幅、分类、版本号）
  \item 签字栏信息（设计者、审核者、工艺、标准化、批准、归档日期）
  \item 材料信息（材质、材料牌号、材料规格）
  \item 滤网规格（编织方式、支数、目数、防菌型）
  \item 性能指标（耐油性、耐酸性、耐热性、耐寒性、抗病毒性能等）
  \item 技术要求（加工要求、表面处理、喷涂颜色、退火处理、拔模斜度）
  \item 部品/BOM 表信息
  \item 其它技术参数（使用温度范围、成形方法、色调、产品名称等）
\end{itemize}

% ==================== 5. NLP 交互与 Agent 调度 ====================
\section{NLP 交互与 Agent 调度}

\subsection{模块功能}

NLP 交互模块（\texttt{nlp\_interface/}）负责解析用户自然语言问题，进行意图分类和参数提取。Agent 核心调度层（\texttt{agent/}）协调各模块工作，实现规则管道与 LLM 回退的混合策略。

\subsection{意图分类体系}

\texttt{intent\_classifier.py} 基于正则表达式将用户问题分为以下意图类型：

\begin{table}[H]
\centering
\begin{tabular}{lll}
\toprule
意图类型 & 示例 & 处理方式 \\
\midrule
\texttt{query\_parameter} & “过滤网的高度是多少？” & 查参数索引→返回值 \\
\texttt{query\_multi} & “这个图纸上有哪些参数？” & 列出所有参数 \\
\texttt{query\_compare} & “外框和滤网哪个更宽？” & 多参数查询+比较 \\
\texttt{query\_info} & “这个零件是做什么的？” & 返回图纸元信息 \\
\bottomrule
\end{tabular}
\end{table}

\subsection{规则管道优先级}

\texttt{query\_processor.py} 实现了严格的规则管道，优先级从高到低依次为：

\begin{enumerate}[nosep]
  \item \textbf{极值查询}（最大/最小/最长/最高/最宽/最窄）
  \item \textbf{计数查询}（几处/几个/多少处）
  \item \textbf{数值搜索}（“5mm 代表什么”）
  \item \textbf{多参数查询}（列表所有参数）
  \item \textbf{比较查询}（A 和 B 哪个更…）
  \item \textbf{单参数查询}（直接匹配参数名）
  \item \textbf{兜底}：列出可查询的参数
\end{enumerate}

\subsection{LLM 回退机制}

当规则管道无法匹配（返回“未找到”或“抱歉”）或检测到复杂问题（如签字栏、标题栏、技术要求等涉及多条件推理的问题）时，系统触发 LLM 回退：

\begin{itemize}[nosep]
  \item 使用 \textbf{DeepSeek v4 Pro}（Anthropic 兼容 API）。
  \item 将全部 102+ 个参数的名称、值、单位、原始文本注入系统提示词。
  \item 同时注入全部文字注释（含图层和位置信息），为 LLM 提供完整的图纸上下文。
  \item 对话历史保留最近 6 轮，支持上下文相关追问。
\end{itemize}

\subsection{文件结构}

\begin{lstlisting}[language={},frame=single]
nlp_interface/               agent/
+-- __init__.py               +-- __init__.py
+-- intent_classifier.py      +-- core.py         # Agent 主循环
+-- parameter_parser.py       +-- session.py      # 会话管理
+-- query_processor.py        +-- config.py       # 配置管理
+-- response_formatter.py     +-- exceptions.py   # 自定义异常
+-- prompt_templates.py       +-- cli.py          # CLI 入口
\end{lstlisting}

% ==================== 6. 考核答案生成 ====================
\section{考核答案生成与评估}

\subsection{考核概述}

竞赛平台发布了两张全新 CAD 图纸作为考核试题：

\begin{itemize}[nosep]
  \item \textbf{侧板图}（侧板MODIFY.dxf）：20 道试题，涉及标题栏、签字栏、材料、技术要求等。
  \item \textbf{印刷版图}（印刷件MODIFY.dxf）：21 道试题，涉及图幅比例、签字栏、材料规格、运输要求等。
\end{itemize}

共计 41 道试题，覆盖尺寸查询、签字栏、标题栏、材料、技术要求、性能参数、实物特征等多类题型。

\subsection{答案生成流程}

\texttt{tools/generate\_exam\_answers.py} 实现自动答案生成：

\begin{enumerate}[nosep]
  \item 读取考核试题 Excel 文件（\texttt{大赛CAD问题.xlsx}）。
  \item 按图纸分组逐图构建独立的参数索引（避免缓存干扰）。
  \item 配置 Agent 启用 LLM 增强，逐个回答问题。
  \item 输出为 \texttt{submission.csv}（含序号、问题、答案）和 \texttt{answers.json}（完整 JSON）。
\end{enumerate}

\subsection{评估方法}

\texttt{tools/evaluate.py} 实现量化评估：

\begin{itemize}[nosep]
  \item 从问题文件中读取 Q$\mid$A 对，逐题调用 Agent 回答。
  \item 数值比较：正则提取答案和预期中的数值，允许 $\pm$0.1 误差。
  \item 文本比较：预期文本是否包含在答案中。
  \item 训练集 58 题评估结果：\textbf{准确率 93.1\%}（从初始 43.1\% 优化而来）。
\end{itemize}

% ==================== 7. 使用与验证 ====================
\section{使用与验证方法}

\subsection{环境要求}

\begin{itemize}[nosep]
  \item Python $\ge$ 3.12
  \item 核心依赖：ezdxf ($\ge$1.4.0)、pydantic ($\ge$2.5.0)、jieba ($\ge$0.42.1)、httpx ($\ge$0.25.0)
  \item 可选依赖：anthropic ($\ge$0.49.0，用于 DeepSeek API 调用)
  \item 转换工具：LibreDWG（\texttt{dwg2dxf}）
  \item 操作系统：Linux（Docker 部署，Ubuntu 24.04）
\end{itemize}

\subsection{快速运行}

\begin{lstlisting}[frame=single,language=bash]
# 1. 安装依赖
pip install -r requirements.txt

# 2. 转换 DWG → DXF（如有需要）
python tools/dwg_to_dxf.py --input data/cad/filter_modify.dwg

# 3. 构建参数索引
python tools/build_index.py \
  --input data/cad/filter_modify.dxf \
  --output data/index/

# 4. 启动交互式问答
python main.py --file data/cad/filter_modify.dxf

# 5. 评估准确率
python tools/evaluate.py \
  --file data/cad/filter_modify.dxf \
  --questions examples/training_questions.txt
\end{lstlisting}

\subsection{考核答案生成}

\begin{lstlisting}[frame=single,language=bash]
python tools/generate_exam_answers.py
# 输出: data/exam_temp/submission.csv
\end{lstlisting}

% ==================== 8. 项目文件结构 ====================
\section{项目代码文件结构}

\begin{lstlisting}[language={},frame=single,title=项目完整目录]
AstraDraft/
+-- main.py                          # 顶层入口
+-- requirements.txt                 # Python 依赖
+-- Dockerfile                       # 开发环境配置
|
+-- agent/                           # Agent 核心调度层
|   +-- core.py                      # 主循环（规则+LLM混合）
|   +-- cli.py                       # CLI 入口
|   +-- config.py                    # 配置管理
|   +-- session.py                   # 对话会话管理
|   +-- exceptions.py                # 自定义异常
|
+-- cad_parser/                      # CAD 解析层
|   +-- reader.py                    # DXF/DWG 文件读取
|   +-- entities.py                  # 几何实体提取
|   +-- dimensions.py                # 尺寸标注提取
|   +-- text_extractor.py            # 文本标注提取
|   +-- layer_manager.py             # 图层管理
|   +-- models.py                    # Pydantic 数据模型
|
+-- cad_understanding/               # 语义理解层
|   +-- parameter_index.py           # 参数索引（核心）
|   +-- semantic_mapper.py           # 语义映射（jieba）
|   +-- component_recognizer.py      # 部件识别
|   +-- scale_inferencer.py          # 单位推断
|
+-- nlp_interface/                   # NLP 交互层
|   +-- intent_classifier.py         # 意图分类
|   +-- parameter_parser.py          # 参数名提取
|   +-- query_processor.py           # 查询处理主流程
|   +-- response_formatter.py        # 回答格式化
|   +-- prompt_templates.py          # LLM 提示词模板
|
+-- tools/                           # 工具脚本
|   +-- dwg_to_dxf.py                # DWG 转换工具
|   +-- build_index.py               # 索引构建
|   +-- evaluate.py                  # 准确率评估
|   +-- generate_exam_answers.py     # 考核答案生成
|
+-- data/
|   +-- cad/                         # CAD 源文件
|   +-- index/                       # 参数索引缓存
|
+-- examples/
    +-- training_questions.txt       # 58 题训练集
\end{lstlisting}

% ==================== 9. 关键技术决策 ====================
\section{关键技术决策与经验总结}

\subsection{结构化解析优于视觉模型}

CAD 图纸本质是结构化数据（矢量图形 + 尺寸标注），而非光栅图片。使用 \textbf{ezdxf} 直接读取结构和标注值，无需多模态/视觉模型。尺寸标注的 \texttt{measurement} 字段直接提供数值，周围文字注释提供语义信息，两者结合即可构建完整参数索引。

\subsection{规则为主、LLM 为辅的混合策略}

\begin{itemize}[nosep]
  \item 规则管道覆盖常见查询类型（尺寸、计数、极值、比较），响应快、结果稳定、无额外成本。
  \item LLM 仅在规则管道无法处理时作为回退，用于复杂推理（如签字栏角色与签名/日期的关联判断、技术要求的综合解读）。
  \item 这种策略兼顾了速度和准确性，并在 LLM API 不可用时保证系统基础功能可用。
\end{itemize}

\subsection{通用化设计}

参数索引构建面向 DXF 标准格式而非特定图纸。标题栏提取、签字栏提取、BOM 表解析都采用基于位置和图层关键词的通用策略，能够适应不同图纸的结构变化。

% ==================== 10. 总结 ====================
\section{总结与展望}

\subsection{关键成果}

\begin{itemize}[nosep]
  \item 成功构建了完整的 CAD 智能问答 Agent 系统，实现了从 DWG/DXF 图纸解析到自然语言问答的全链路。
  \item 训练集 58 题评估准确率 \textbf{93.1\%}，支持 12 类以上查询意图。
  \item 参数索引覆盖 108+ 个参数，新增编织方式、防菌型、退火处理、滤网支数、部品表等提取规则。
	  \item 自动生成了 41 道考核试题的答案（侧板 20 题 + 印刷版 21 题）。
  \item 参赛代码和本技术方案报告已按要求准备完毕。
\end{itemize}

\subsection{改进方向}

\begin{itemize}[nosep]
  \item 进一步优化标题栏和签字栏的关联识别精度，特别是在不同图纸布局下的自适应能力。
  \item 扩展对更多 CAD 实体类型和标注类型的支持，提高参数覆盖率。
  \item 探索使用更多 LLM 能力进行多轮对话和复杂推理，同时保持规则管道的稳定性。
\end{itemize}

\end{document}
